\documentclass[12pt,a4paper]{article}

\usepackage{cmap}
\usepackage[T2A]{fontenc}
\usepackage[utf8]{inputenc}
\usepackage[russian]{babel}
\usepackage{amsmath,amssymb}
\usepackage[a4paper,margin=2cm]{geometry}

\newcommand{\R}{\mathbb{R}}
\newcommand{\eps}{\varepsilon}
\newcommand{\dd}{\,d}
\newcommand{\grad}{\operatorname{grad}}
\newcommand{\Int}{\operatorname{int}}
\newcommand{\Jac}{D}

\setlength{\parindent}{0pt}
\setlength{\parskip}{0.45em}
\sloppy

\begin{document}

\begin{center}
{\Large\bfseries Билет 3}
\end{center}

\noindent\fbox{%
  \begin{minipage}{\dimexpr\linewidth-2\fboxsep-2\fboxrule\relax}
  \normalfont
Частные производные функции нескольких переменных. Дифференцируемость
функции в точке, дифференциал. Необходимые условия дифференцируемости,
достаточные условия дифференцируемости функции нескольких переменных.
Дифференцируемость сложной функции. Инвариантность формы дифференциала
относительно замены переменных. Градиент и производная по направлению.
  \end{minipage}%
}

\section*{Краткий план}

Сначала определяется дифференцируемость функции \(f:\R^n\to\R\) как наличие
главной линейной части приращения. Эта линейная часть задается градиентом, а
сам дифференциал равен скалярному произведению градиента на приращение
аргумента. Затем вводятся производные по направлению и частные производные.
Главные факты: дифференцируемость влечет непрерывность, существование всех
производных по направлениям и всех частных производных; при непрерывности
частных производных в точке функция дифференцируема. Для вектор-функций
дифференциал записывается через матрицу Якоби, а для композиции работает
матричное правило цепочки. Из него следует инвариантность формы первого
дифференциала.

\section*{Определения}

\textbf{Дифференцируемость скалярной функции.}
Пусть функция \(f(x)=f(x_1,\dots,x_n)\) определена в окрестности точки
\(x^0=(x^0_1,\dots,x^0_n)\in\R^n\). Функция \(f\) называется
дифференцируемой в точке \(x^0\), если существует вектор
\(A=(A_1,\dots,A_n)\in\R^n\) такой, что при \(x\to x^0\)
\[
  f(x)-f(x^0)=(A,x-x^0)+o(|x-x^0|).
\]
Этот вектор называется градиентом функции в точке \(x^0\) и обозначается
\[
  A=\grad f(x^0).
\]
Значит, для дифференцируемой функции
\[
  f(x)-f(x^0)=(\grad f(x^0),x-x^0)+o(|x-x^0|).
\]

\textbf{Дифференциал.}
Дифференциалом функции \(f\) в точке \(x^0\) называется линейная относительно
приращения \(dx=x-x^0\) функция
\[
  df(x^0)=(\grad f(x^0),dx).
\]
Если \(f\) дифференцируема в \(x^0\), то
\[
  \Delta f=f(x)-f(x^0)=df(x^0)+o(|dx|).
\]
После установления существования частных производных дифференциал записывается
в координатной форме:
\[
  df(x^0)=\sum_{i=1}^n \frac{\partial f}{\partial x_i}(x^0)\,dx_i,
  \qquad dx_i=x_i-x_i^0.
\]

\textbf{Производная по вектору и по направлению.}
Производной функции \(f\) в точке \(x^0\) по вектору \(\ell\in\R^n\)
называется правый предел
\[
  \frac{\partial f}{\partial \ell}(x^0)
  =
  \lim_{t\to+0}\frac{f(x^0+t\ell)-f(x^0)}{t},
\]
если он существует. Если \(|\ell|=1\), то \(\ell\) называется направлением, а
эта производная называется производной по направлению.

\textbf{Частная производная.}
Частной производной функции \(f(x_1,\dots,x_n)\) по переменной \(x_i\) в точке
\(x^0\) называется обычная производная по \(x_i\) при фиксированных остальных
переменных:
\[
  \frac{\partial f}{\partial x_i}(x^0)
  =
  \lim_{x_i\to x_i^0}
  \frac{
    f(x_1^0,\dots,x_{i-1}^0,x_i,x_{i+1}^0,\dots,x_n^0)-f(x^0)
  }{x_i-x_i^0}.
\]

\textbf{Вектор-функция и матрица Якоби.}
Пусть \(F=(F_1,\dots,F_m):\R^n\to\R^m\). Вектор-функция \(F\) называется
дифференцируемой в точке \(x^0\), если все ее координаты \(F_k\)
дифференцируемы в этой точке. Ее дифференциал -- вектор из дифференциалов
координат:
\[
  dF(x^0)=
  \begin{pmatrix}
    dF_1(x^0)\\
    \vdots\\
    dF_m(x^0)
  \end{pmatrix}.
\]
Матрицей Якоби \(F\) в точке \(x^0\) называется матрица
\[
  \Jac F(x^0)=
  \begin{pmatrix}
    \dfrac{\partial F_1}{\partial x_1}(x^0) & \cdots &
    \dfrac{\partial F_1}{\partial x_n}(x^0)\\
    \vdots & \ddots & \vdots\\
    \dfrac{\partial F_m}{\partial x_1}(x^0) & \cdots &
    \dfrac{\partial F_m}{\partial x_n}(x^0)
  \end{pmatrix}.
\]
В матричной форме
\[
  dF(x^0)=\Jac F(x^0)\,dx.
\]

\section*{Теоремы}

\textbf{1. Необходимые условия дифференцируемости.}
Пусть \(f\) определена в окрестности \(x^0\in\R^n\) и дифференцируема в этой
точке. Тогда:
\[
  f \text{ непрерывна в } x^0;
\]
для любого \(\ell\in\R^n\) существует производная по вектору \(\ell\), причем
\[
  \frac{\partial f}{\partial \ell}(x^0)
  =
  (\grad f(x^0),\ell);
\]
все частные производные существуют и
\[
  \grad f(x^0)=
  \left(
    \frac{\partial f}{\partial x_1}(x^0),\dots,
    \frac{\partial f}{\partial x_n}(x^0)
  \right).
\]
Следовательно,
\[
  df(x^0)=\sum_{i=1}^n\frac{\partial f}{\partial x_i}(x^0)\,dx_i.
\]

\textbf{2. Достаточное условие дифференцируемости.}
Если все частные производные
\[
  \frac{\partial f}{\partial x_i},\qquad i=1,\dots,n,
\]
определены в некоторой окрестности точки \(x^0\in\R^n\) и непрерывны в точке
\(x^0\), то \(f\) дифференцируема в \(x^0\).

\textbf{3. Дифференцирование сложной вектор-функции.}
Пусть \(X\subset\R^n\), \(Y\subset\R^m\), \(F:X\to Y\), \(G:Y\to\R^p\).
Пусть \(F\) дифференцируема в точке \(x^0\in\Int X\), а \(G\)
дифференцируема в точке \(y^0=F(x^0)\in\Int Y\). Тогда
\[
  \Phi(x)=G(F(x))
\]
дифференцируема в точке \(x^0\), и
\[
  \Jac\Phi(x^0)=\Jac G(y^0)\,\Jac F(x^0).
\]
В координатах:
\[
  \frac{\partial \Phi_i}{\partial x_j}(x^0)
  =
  \sum_{k=1}^m
  \frac{\partial G_i}{\partial y_k}(y^0)
  \frac{\partial F_k}{\partial x_j}(x^0),
  \qquad i=1,\dots,p,\quad j=1,\dots,n.
\]

\textbf{4. Инвариантность формы первого дифференциала.}
Пусть \(y=F(x)\) дифференцируема в \(x^0\in\R^n\), а \(z=G(y)\)
дифференцируема в \(y^0=F(x^0)\in\R^m\). Тогда для простой функции
\(z=G(y)\) и для сложной функции \(z=\Phi(x)=G(F(x))\) формула первого
дифференциала имеет один и тот же вид:
\[
  dz=\Jac G(y^0)\,dy.
\]
В случае простой функции \(dy\) -- приращение независимой переменной \(y\), а
в случае сложной функции \(dy=dF(x^0)=\Jac F(x^0)\,dx\).

\textbf{5. Геометрический смысл градиента.}
Если \(f\) дифференцируема в \(x^0\) и \(\grad f(x^0)\ne0\), то среди всех
единичных направлений производная по направлению максимальна в направлении
\[
  \ell_{\max}=\frac{\grad f(x^0)}{|\grad f(x^0)|}
\]
и минимальна в направлении
\[
  \ell_{\min}=-\frac{\grad f(x^0)}{|\grad f(x^0)|}.
\]
Максимальное значение равно \(|\grad f(x^0)|\), минимальное равно
\(-|\grad f(x^0)|\).

\section*{Доказательства}

\textbf{Необходимые условия.}
Пусть \(f\) дифференцируема в \(x^0\). Тогда
\[
  f(x)-f(x^0)=(\grad f(x^0),x-x^0)+o(|x-x^0|).
\]
При \(x\to x^0\) линейная часть стремится к нулю, а остаток
\(o(|x-x^0|)\) также стремится к нулю. Значит
\[
  f(x)\to f(x^0),
\]
то есть \(f\) непрерывна в \(x^0\).

Зафиксируем \(\ell\in\R^n\) и положим \(x=x^0+t\ell\), где \(t\to+0\). Тогда
\[
  f(x^0+t\ell)-f(x^0)
  =
  (\grad f(x^0),t\ell)+o(t),
\]
потому что \(|t\ell|=O(t)\). Делим на \(t\) и получаем
\[
  \frac{\partial f}{\partial \ell}(x^0)
  =
  \lim_{t\to+0}
  \frac{f(x^0+t\ell)-f(x^0)}{t}
  =
  (\grad f(x^0),\ell).
\]

Возьмем координатные направления \(e_i=(0,\dots,0,1,0,\dots,0)\) и
\(-e_i\). По только что доказанной формуле существуют односторонние
производные вдоль этих направлений:
\[
  \frac{\partial f}{\partial e_i}(x^0)=(\grad f(x^0),e_i),
  \qquad
  \frac{\partial f}{\partial (-e_i)}(x^0)=-(\grad f(x^0),e_i).
\]
Правая и левая производные функции одной переменной, полученной фиксацией всех
переменных кроме \(x_i\), существуют и согласованы; поэтому существует частная
производная по \(x_i\), причем она равна \(i\)-й координате градиента:
\[
  \frac{\partial f}{\partial x_i}(x^0)=(\grad f(x^0))_i.
\]
Отсюда сразу следует координатная формула для градиента и дифференциала:
\[
  df(x^0)=(\grad f(x^0),dx)
  =
  \sum_{i=1}^n
  \frac{\partial f}{\partial x_i}(x^0)\,dx_i.
\]

\textbf{Достаточное условие.}
Обозначим \(h=x-x^0\), \(x=x^0+h\). Приращение функции можно разложить по
координатам. Для этого положим
\[
  y^{(i)}=(x^0_1+h_1,\dots,x^0_i+h_i,x^0_{i+1},\dots,x^0_n),
  \qquad y^{(0)}=x^0.
\]
Тогда
\[
  f(x^0+h)-f(x^0)
  =
  \sum_{i=1}^n \bigl(f(y^{(i)})-f(y^{(i-1)})\bigr).
\]
К каждому слагаемому применяем одномерную теорему Лагранжа по переменной
\(x_i\). Получаем точки \(\xi_i\), лежащие между соответствующими концами
отрезков, такие что
\[
  f(x^0+h)-f(x^0)
  =
  \sum_{i=1}^n
  \frac{\partial f}{\partial x_i}(\xi_i)\,h_i.
\]
Так как \(\xi_i\to x^0\) при \(h\to0\) и частные производные непрерывны в
\(x^0\),
\[
  \frac{\partial f}{\partial x_i}(\xi_i)
  =
  \frac{\partial f}{\partial x_i}(x^0)+\eps_i(h),
  \qquad \eps_i(h)\to0.
\]
Следовательно,
\[
  f(x^0+h)-f(x^0)
  =
  \sum_{i=1}^n
  \frac{\partial f}{\partial x_i}(x^0)h_i
  +
  \sum_{i=1}^n \eps_i(h)h_i.
\]
Остаток является \(o(|h|)\), поскольку
\[
  \left|\sum_{i=1}^n\eps_i(h)h_i\right|
  \le
  \max_i|\eps_i(h)|\sum_{i=1}^n |h_i|
  \le
  \sqrt n\,\max_i|\eps_i(h)|\,|h|
  =
  o(|h|).
\]
Значит
\[
  f(x)-f(x^0)=
  \sum_{i=1}^n
  \frac{\partial f}{\partial x_i}(x^0)(x_i-x_i^0)
  +
  o(|x-x^0|),
\]
то есть \(f\) дифференцируема в \(x^0\).

\textbf{Дифференцирование сложной функции.}
Пусть \(h=x-x^0\). По дифференцируемости \(F\) и \(G\)
\[
  F(x)-F(x^0)=\Jac F(x^0)h+o(|h|),
\]
\[
  G(y)-G(y^0)=\Jac G(y^0)(y-y^0)+o(|y-y^0|).
\]
Подставим \(y=F(x)\), \(y^0=F(x^0)\). Так как дифференцируемость \(F\)
влечет непрерывность и даже оценку \(|F(x)-F(x^0)|=O(|h|)\), имеем
\[
  o(|F(x)-F(x^0)|)=o(|h|).
\]
Поэтому
\[
  G(F(x))-G(F(x^0))
  =
  \Jac G(y^0)\Jac F(x^0)h+o(|h|).
\]
По определению дифференцируемости вектор-функции это означает, что
\(\Phi=G\circ F\) дифференцируема в \(x^0\), а ее матрица Якоби равна
\[
  \Jac\Phi(x^0)=\Jac G(y^0)\Jac F(x^0).
\]
Координатная формула получается из правила умножения матриц.

\textbf{Инвариантность формы дифференциала.}
Для простой функции \(z=G(y)\) по формуле дифференциала вектор-функции
\[
  dz=dG(y^0)=\Jac G(y^0)\,dy.
\]
Для сложной функции \(z=\Phi(x)=G(F(x))\)
\[
  dz=d\Phi(x^0)=\Jac\Phi(x^0)\,dx.
\]
По правилу цепочки
\[
  \Jac\Phi(x^0)=\Jac G(y^0)\Jac F(x^0),
\]
а
\[
  dy=dF(x^0)=\Jac F(x^0)\,dx.
\]
Следовательно,
\[
  dz=\Jac G(y^0)\Jac F(x^0)\,dx=\Jac G(y^0)\,dy.
\]
Форма записи такая же, как у простой функции \(z=G(y)\), только в случае
замены переменных \(dy\) уже означает дифференциал функции \(y=F(x)\).

\textbf{Градиент и направление наибольшего роста.}
Для единичного направления \(|\ell|=1\) по необходимому условию
\[
  \frac{\partial f}{\partial\ell}(x^0)=(\grad f(x^0),\ell).
\]
По неравенству Коши--Буняковского
\[
  (\grad f(x^0),\ell)\le |\grad f(x^0)|\,|\ell|=|\grad f(x^0)|.
\]
Равенство достигается при
\[
  \ell=\frac{\grad f(x^0)}{|\grad f(x^0)|}.
\]
Аналогично минимум достигается в противоположном направлении
\[
  \ell=-\frac{\grad f(x^0)}{|\grad f(x^0)|},
\]
и равен \(-|\grad f(x^0)|\).

\section*{Финальная устная версия}

Функция \(f:\R^n\to\R\) дифференцируема в точке \(x^0\), если ее приращение
имеет главную линейную часть:
\[
  f(x)-f(x^0)=(A,x-x^0)+o(|x-x^0|).
\]
Вектор \(A\) называется градиентом и обозначается \(\grad f(x^0)\), а
дифференциал равен
\[
  df(x^0)=(\grad f(x^0),dx).
\]
Если функция дифференцируема, то она непрерывна в этой точке, потому что и
линейная часть, и остаток стремятся к нулю.

Производная по вектору \(\ell\) определяется как
\[
  \frac{\partial f}{\partial\ell}(x^0)=
  \lim_{t\to+0}\frac{f(x^0+t\ell)-f(x^0)}{t}.
\]
Для дифференцируемой функции она существует для любого \(\ell\) и равна
\[
  \frac{\partial f}{\partial\ell}(x^0)=(\grad f(x^0),\ell).
\]
Если \(|\ell|=1\), это производная по направлению. Частная производная по
\(x_i\) -- это обычная производная по \(x_i\) при фиксированных остальных
переменных. Из дифференцируемости следуют существование всех частных
производных и формула
\[
  \grad f(x^0)=
  \left(
    \frac{\partial f}{\partial x_1}(x^0),\dots,
    \frac{\partial f}{\partial x_n}(x^0)
  \right),
  \qquad
  df(x^0)=\sum_{i=1}^n\frac{\partial f}{\partial x_i}(x^0)\,dx_i.
\]

Существование частных производных само по себе еще не гарантирует
дифференцируемость. Достаточное условие такое: если все частные производные
определены в окрестности точки и непрерывны в самой точке, то функция
дифференцируема. Доказательство получается разложением приращения по
координатам и применением одномерной теоремы Лагранжа к каждому слагаемому;
непрерывность частных производных дает остаток \(o(|x-x^0|)\).

Для вектор-функции \(F:\R^n\to\R^m\) дифференцируемость означает
дифференцируемость всех координат, а дифференциал записывается через матрицу
Якоби:
\[
  dF(x^0)=\Jac F(x^0)\,dx.
\]
Если \(F\) дифференцируема в \(x^0\), а \(G\) дифференцируема в
\(y^0=F(x^0)\), то сложная функция \(\Phi=G\circ F\) дифференцируема и
\[
  \Jac\Phi(x^0)=\Jac G(y^0)\Jac F(x^0).
\]
В координатах это правило цепочки:
\[
  \frac{\partial \Phi_i}{\partial x_j}(x^0)
  =
  \sum_{k=1}^m
  \frac{\partial G_i}{\partial y_k}(y^0)
  \frac{\partial F_k}{\partial x_j}(x^0).
\]

Инвариантность формы первого дифференциала означает, что для \(z=G(y)\) и при
замене \(y=F(x)\) остается одна и та же запись
\[
  dz=\Jac G(y^0)\,dy.
\]
Разница только в смысле \(dy\): для простой функции это приращение независимой
переменной, а для сложной функции это дифференциал \(dF(x^0)\).

Наконец, если \(\grad f(x^0)\ne0\), то направление градиента является
направлением наибольшего роста. Действительно, для \(|\ell|=1\)
\[
  \frac{\partial f}{\partial\ell}(x^0)=(\grad f(x^0),\ell)
  \le |\grad f(x^0)|.
\]
Равенство достигается при
\(\ell=\grad f(x^0)/|\grad f(x^0)|\), а минимум -- в противоположном
направлении.

\section*{Источники}

Иванов Г.Е. \emph{Лекции по математическому анализу. Часть 1}:
\texttt{IvGE\_1.pdf}, стр. 180--191.

Основные роли источника: стр. 180--183 -- дифференцируемость, градиент,
дифференциал и геометрический смысл; стр. 183--186 -- необходимые условия,
производные по направлениям и частные производные; стр. 187--188 --
достаточное условие дифференцируемости; стр. 188--191 -- матрица Якоби,
дифференцирование сложной вектор-функции и инвариантность формы первого
дифференциала.

\end{document}
